% Generated by Listas/make_statements.py from pratica-lecture-05.md, item 10. Do not edit.
\textbf{Source.} Tadelis (2013), \emph{Game Theory: An Introduction}, Exercise 10.3 ``Not-So-Grim Trigger'', pp.~215--216 {[}PDF 234--235{]}, with its footnote 12 (p.~216). Answer in the source: none (the book prints no solutions; the preface says about half are on Princeton UP's website, not checked). Reproduced verbatim. (Here \(T\) is the length of the punishment, not a payoff.)



\begin{quote}

\textbf{Not-So-Grim Trigger:} Consider the infinitely repeated Prisoner's Dilemma with discount factor \(\delta < 1\) described by the following matrix:



{\def\LTcaptype{none} % do not increment counter

\begin{longtable}[]{@{}ccc@{}}

\toprule\noalign{}

Player 1 ~Player 2 & \(m\) & \(f\) \\

\midrule\noalign{}

\endhead

\bottomrule\noalign{}

\endlastfoot

\(M\) & 4, 4 & −1, 5 \\

\(F\) & 5, −1 & 1, 1 \\

\end{longtable}

}



Instead of using grim-trigger strategies to support a pair of actions \((a_1, a_2)\) other than \((F, f)\) as a subgame-perfect equilibrium, assume that the players wish to choose a less draconian punishment called a ``length-\(T\) punishment'' strategy. If there is a deviation from \((a_1, a_2)\) then the players will play \((F, f)\) for \(T\) periods and then resume playing \((a_1, a_2)\). Let \(\delta_T\) be the critical discount factor so that if \(\delta > \delta_T\) then the adequately defined strategies will implement the desired path of play with length-\(T\) punishment as the threat.



\begin{enumerate}

\def\labelenumi{\alph{enumi}.}

\item

  Let \(T = 1\). What is the critical value \(\delta_1\) to support the pair of actions \((M, m)\) played in every period?

\item

  Let \(T = 2\). What is the critical value \(\delta_T\) to support the pair of actions \((M, m)\) played in every period?¹²

\item

  Compare the two critical values in (a) and (b). How do they differ and what is the intuition for this?

\end{enumerate}



¹² Hint: You should encounter an equation of the form \(a\delta^3 - (a + 1)\delta + 1 = 0\), for which it is easy to see that \(\delta = 1\) is a root. In this case you know that the equation can be written in the form \((\delta - 1)(a\delta^2 + a\delta - 1) = 0\), and you can solve for the other relevant root of the cubic equation.

\end{quote}

